% Schematic drawing of one possible shape of the separation graph.
\begin{figure}[!htbp]
  \centering
  \begin{tikzpicture}[x=1cm,y=1cm,font=\normalsize,
      axis/.style={line width=.55pt,-{Stealth[length=2mm]}},
      graph/.style={line width=.9pt},
      guide/.style={line width=.45pt,densely dashed},
      projection/.style={line width=.4pt,densely dotted}]
    \node at (5.1,3.0)
      {$\Delta(r,m,p)=r^2V(r,m,p)=r^2\bigl[mH_\Delta(r,m,p)+rB_\Delta(r,p)\bigr]$};

    \draw[axis] (0,-1.15) -- (0,2.55)
      node[above] {$\delta$};
    \node[anchor=east,align=right] at (-.22,2.12)
      {threshold\\separation};
    \draw[axis] (0,0) -- (10.05,0)
      node[below] {$r=\sqrt q$};
    \node[below left] at (0,0) {$0$};

    % The two selected intersections lie on this hand-drawn guide.
    \draw[graph]
      (0,0)
      .. controls (.8,0) and (1.55,.45) .. (2.65,1.1)
      .. controls (3.75,1.75) and (4.55,2.05) .. (5.35,2.05)
      .. controls (6.05,2.05) and (7.1,1.72) .. (7.8,1.1)
      .. controls (8.45,.52) and (9.0,-.32) .. (9.32,-.94);
    \node[anchor=east] at (8.5,-.94) {$\Delta(r,m,p)$};

    \draw[guide] (0,1.1) -- (9.5,1.1)
      node[anchor=west] {$\delta=\text{constant}$};
    \draw[projection] (2.65,0) -- (2.65,1.1);
    \draw[projection] (7.8,0) -- (7.8,1.1);
    \draw[projection] (5.35,0) -- (5.35,2.05);
    \fill (2.65,1.1) circle[radius=2pt];
    \fill (7.8,1.1) circle[radius=2pt];
    \node[below] at (2.65,0) {$r_1$};
    \node[below] at (7.8,0) {$r_2$};
    \node[below] at (5.35,0) {$r_{\rm c}$};
    \node[above] at (5.35,2.08) {positive critical point};
  \end{tikzpicture}
  \caption{The separation graph $\delta=\Delta(r,m,p)$ for fixed $(m,p)$ with
    $c_0<0$, $A>0$ and $m>0$ (schematic). The graph has at most one positive
    critical point, which is nondegenerate
    (Proposition~\ref{prop:one-fold-graph}). Its intersections with a
    horizontal line are the zeros of $P''$ in the outer region $q>L|\delta|$,
    by~\eqref{eq:curvature-factor}; for other parameters the graph can be
    monotone. The resulting bound on fixed points is
    Proposition~\ref{prop:common-count}.}
  \label{fig:curvature-graph}
\end{figure}
